\vspace{-0.2 cm}
\section{Policy Decorator: Model-Agnostic Online Refinement}
\vspace{-0.2 cm}

In this work, our goal is to online improve a large policy model, which is usually offline-trained by imitation learning and usually has some specific designs in model architecture. To this end, we propose \textit{Policy Decorator}, a model-agnostic framework for refining large policy models via online interactions with environments. Fig. \ref{fig:method_overview} provides an overview of our framework.

Policy Decorator is grounded on \textbf{learning a residual policy via reinforcement learning with sparse rewards}, which is described in Sec. \ref{sec:residual_rl}. On top of it, we devise a set of strategies to ensure the RL agent (in combination with the base policy and the residual policy) explores the environment in a controlled manner. Such a \textbf{controlled exploration} mechanism is detailed in Sec. \ref{sec:control_explore}. Finally, we discuss several important \textbf{design choices} that further enhance learning efficiency in Sec. \ref{sec:design_choices}.



\subsection{Learning Residual Policy via RL}
\label{sec:residual_rl}

Given the base policy $\pi_{base}$, we then train a residual policy $\pi_{res}$ on top of it using reinforcement learning. The base policy $\pi_{base}$ can be either deterministic (e.g., Behavior Transformer \citep{bet}) or stochastic (e.g., Diffusion Policy \citep{dpi}), and it remains frozen during the RL process. The residual policy $\pi_{res}$ is updated through RL gradients, so it should be a differentiable function compatible with RL gradients. In this work, we model the residual policy $\pi_{res}$ as a Gaussian policy parameterized by a small neural network (either an MLP or a CNN, depending on the observation modality).
To interact with the environment, the actions from both policies are combined by summing their output actions, i.e., the action executed in the environment is $\pi_{base}(s) + \pi_{res}(s)$. For stochastic policies, actions are sampled individually from both policies and then summed.

The residual policy is trained to maximize the expected discounted return derived from the sparse reward (i.e., the task's success signal). We employ the Soft Actor-Critic (SAC) algorithm \citep{sac} due to its superior sample efficiency and stability. Several important design choices arise when implementing SAC for learning the residual policy, which we discuss in Sec. \ref{sec:design_choices}.
Our method is also compatible with PPO \citep{ppo}, as illustrated in Appendix \ref{appendix:ppo}. 


\subsection{Controlled Exploration}
\label{sec:control_explore}

While learning a residual policy by RL can in principle refine a base policy, practical implementation can be challenging. 
As demonstrated in our experiments (Sec. \ref{sec:experiment}), without constraints, random exploration during RL training often leads to failure in tasks requiring precise control, resulting in no learning signals in sparse reward settings (see \href{https://policydecorator.github.io/random_residual_actions.html}{\color{green} this video} for an example).
To overcome these challenges and enable stable, sample-efficient learning, we propose a set of strategies ensuring the RL agent (in combination with the base policy and the residual policy) \textbf{explores the environment in a controlled manner}. The goal is to make sure the agent continuously receives sufficient success signals while adequately exploring the environment.


\textbf{Bounded Residual Action}~ 
When using the residual policy to correct the base policy, we do not want the resulting trajectory to deviate too much from the original trajectory because it usually leads to failure. Instead, we expect the residual policy to only make a bit "refinement" at the finer parts of the tasks. 
To reflect this spirit, we bound the output of the residual policy within a certain scale. Since we use SAC as our backbone RL algorithm, the output of the policy is naturally bounded by a squashing function (\texttt{tanh}), whose range is $(-1, 1)$. We further scale the action sampled from the Gaussian policy with a hyperparameter $\alpha$, making the range of the residual action $(-\alpha, \alpha)$.
We found that an appropriate scale of residual action bound can be crucial for some precise tasks. We investigated the effects of hyperparameter $\alpha$ in Sec. \ref{sec:ablation_hp}.



\textbf{Progressive Exploration Schedule}~
Given that our residual policy is randomly initialized, the agent (combined with the base policy and residual policy) may exhibit highly random behavior and fail to succeed at the initial stage of learning. Therefore, the base policy alone, trained by imitation learning, can be more reliable during the early stages. As training progresses, the residual policy can be gradually improved, making it safer to incorporate its suggestions.

\begin{wrapfigure}{r}{0.28\linewidth}
    \centering
    \vspace{-0.7cm}
    \includegraphics[width=\linewidth]{figures/offset_schedule.pdf}
    \vspace{-0.8cm}
    \caption{Progressive Exploration Schedule.}
    \label{fig:offset_schedule}
    \vspace{-1cm}
\end{wrapfigure}


Inspired by the $\epsilon$-greedy strategy used in DQN \citep{dqn}, we propose to progressively introduce actions from the residual policy into the agent's behavior policy. Specifically, the behavior policy will use actions from the residual policy to complement the base policy with probability $\epsilon$ and rely solely on the base policy with probability $1-\epsilon$. Formally, during training,
\begin{equation}
    \pi_{behavior}(s) = 
    \begin{cases}
        ~\pi_{base}(s) + \pi_{res}(s) & \text{Uniform}(0,1) < \epsilon \\
        ~\pi_{base}(s) & \text{otherwise}
    \end{cases}
\end{equation}

The parameter $\epsilon$ increases linearly from 0 to 1 over a specified number of time steps, as shown in Fig. \ref{fig:offset_schedule}, where $H$ is a hyperparameter. Our experiments in Sec. \ref{sec:ablation_hp} indicate that while tuning $H$ can enhance sample efficiency, using a large $H$ is generally a safe choice.

\subsection{Design Choices \& Implementation Details}
\label{sec:design_choices}

We investigated a few important design choices in the implementation of Policy Decorator, with supporting experiments provided in Appendix \ref{appendix:design_choices}.

\textbf{Input of Residual Policy}~
The residual policy can receive input in the form of either observation alone or both observation and action from the base policy. Our experiments indicate that using only observation typically produces better results, as illustrated in Fig. \ref{fig:input_of_actor}.

\textbf{Input of Critic}~
In SAC, the critic $Q(s, a)$ takes an action as input, and there are several design choices regarding this action: we can use 1) the sum of the base action and residual action; 2) the concatenation of both; or 3) the residual action alone. Based on our experiments shown in Fig. \ref{fig:input_of_critic}, using the sum of both actions yields the best performance. 










    
    








